\documentclass[12pt, a4paper,oneside]{report}
\usepackage[spanish]{babel}
\selectlanguage{spanish}
\usepackage{multicol}
\usepackage{caption}
\usepackage{graphicx} % Required for inserting images
\usepackage{subcaption}
\usepackage{ragged2e}
\usepackage{color}
\usepackage[dvipsnames,table]{xcolor}
\usepackage{float}
\usepackage{caption}
\captionsetup[table]{skip=10pt}
\usepackage{array,tabularx, longtable}
\usepackage{bigstrut}
\usepackage{multirow}
\usepackage{enumerate} 
\usepackage{enumitem}
\usepackage[T1]{fontenc}
\usepackage{xspace}
\usepackage{newtxtext} % letra Times New Roman
\usepackage{titlesec}


\usepackage{anysize}

\usepackage{blindtext}
%Matematica
\usepackage{times}
\usepackage[utf8]{inputenc}
\usepackage{amsmath}
\usepackage{calc}

%Agregar codigo
\usepackage{verbatim}
\usepackage{listings}

\usepackage{setspace}
\usepackage{booktabs} 

\usepackage[hidelinks]{hyperref}

%\usepackage{lscape} % cambia la direccion
\usepackage{pdflscape}

\usepackage[refpages]{gloss}
\usepackage{glossaries}
\usepackage{acronym}

\usepackage{fancyhdr}%encabezados y pie de pagina  \pagestyle{fancy}


\graphicspath{{imagenes/}}
\usepackage{lipsum}  

\usepackage{appendix} %Anexos
\usepackage{tocloft} %personalizar titulos
\renewcommand{\cftchapleader}{\cftdotfill{\cftdotsep}}

%bibliografia y referencias
\usepackage{apacite}
\usepackage[sectionbib]{natbib}
\usepackage{chapterbib}
\usepackage[hidelinks]{hyperref}



\bibliographystyle{apacite}


%Cambiar tamano tittulos
\titleformat{\chapter}[display]{\normalfont\sffamily\large\bfseries}
{\chaptertitlename\ \thechapter}{12pt}{\large}

\titleformat{\section}
{\normalfont\sffamily\large\bfseries}
{\thesection}{1em}{}

\addtocounter{secnumdepth}{4} % 3 => section, subsection, subsubsection
\renewcommand\thesubsubsection{\arabic{chapter}.\arabic{section}.\arabic{subsection}.\arabic{subsubsection}}
 %FORMATOS APA EN IMAGENES Y TABLAS

\captionsetup[figure]{labelfont=bf, textfont={it, normalsize}, labelsep=newline , justification=raggedright, singlelinecheck=false}
\captionsetup[table]{labelfont=bf, textfont={it, normalsize}, labelsep=newline , justification=raggedright, singlelinecheck=false}
\begin{document}
	\renewcommand{\BOthers}[1]{et al.\hbox{}}

	%=====================MARGENES=============================
	\marginsize{2.54cm}{2.54cm}{2.54cm}{2.54cm}

        %=========== Titulos indices ===========================
	\renewcommand*{\contentsname}{Tabla de contenidos}
	\renewcommand*{\listfigurename}{Índice de figuras}
	\renewcommand*{\listtablename}{Índice de tablas}

	%============ personalizar tipo de numeracion y anexos=====
	\renewcommand{\baselinestretch}{1.0}
	\renewcommand{\appendixname}{Anexos}
	\renewcommand{\appendixtocname}{Anexos}
	\renewcommand{\appendixpagename}{Anexos}
	\renewcommand{\thetable}{\arabic{table}}
	\renewcommand*{\tablename}{Tabla}
	\renewcommand*{\chaptername}{Capítulo}
	\renewcommand*{\thechapter}{\Roman{chapter}}
	\renewcommand{\thesection}{\arabic{chapter}.\arabic{section}}
	\renewcommand{\figurename}{Figura}
	%\renewcommand{\thefigure}{\arabic{chapter}.\arabic{figure}}
    \renewcommand{\thefigure}{\arabic{figure}}
	\renewcommand{\theequation}{\arabic{chapter}.\arabic{equation}}
	\renewcommand\bibname{Bibliografía}
 
\newcommand*{\figuretitle}[1]{%
    {\centering%   <--------  will only affect the title because of the grouping (by the
    \textbf{#1}%              braces before \centering and behind \medskip). If you remove
    \par\medskip}%            these braces the whole body of a {figure} env will be centered.
}
 %==================================================

%============  PORTADA(CARATULA)  ==============
\input{Contenido/Caratula}

%============ INFORME DE SIMILITUD ==============
\input{Contenido/Informe de Similitud}

%========Numeracion de pagina
\pagenumbering{roman}

% =====poner numeracion en el encabezado
\fancypagestyle{plain}{
	\fancyhf{}
	\fancyhead[R]{\thepage}
	\renewcommand{\headrulewidth}{0pt}
	\renewcommand{\footrulewidth}{0pt}
}
%=============== DEDICATORIA ============
\input{Contenido/Dedicatoria}

%=============== AGRADECIMIENTO ============
\input{Contenido/Agradecimiento}

%=================== INDICE GENERAL =====================
\rm
\singlespacing 
\renewcommand\contentsname{\centering \Large Índice General }
%indices previos
\addcontentsline{tocbib}{chapter}{\contentsname}
%=====Tabla de contenido
\large{
\tableofcontents}
\newpage
%=================== LISTA DE CUADROS =====================
\listoftables
\newpage
%=================== INDICE FIGURAS =====================
\listoffigures
\newpage

%Resumen en español y en el idioma acreditado
%============ RESUMEN ===============
\input{Contenido/Resumen}
%============ Palabras clave ============================
\input{Contenido/Palabras Clave}

\pagenumbering{arabic}

%====================CONTENIDO=================
\linespread{2}
%===================== INTRODUCCIÓN =======================

%======== PLANTEAMIENTO DEL PROBLEMA========================
\pagestyle{plain}
\rhead[]{\thepage}
\input{Contenido/Cap I Planteamiento}

%=========== MARCO TEORICO==================
\input{Contenido/Cap II Marco Teorico}

%=========== METODOLOGIA=======================
\input{Contenido/Cap III Metodologia}

%=========== RESULTADOS Y DISCUSIÓN =================
\input{Contenido/Cap IV Resultados y Discusion}

%=======================CONCLUSIONES Y RECOMENDACIONES=======
\input{Contenido/Cap V Conclusiones y Recomendaciones}
\newpage

%=========REFERENCIAS BIBLIOGRÁFICAS=========================

%\renewcommand\contentsname{\centering \large Índice General }
\renewcommand{\bibname}{ \large Referencias Bibliograficas}

\bibliography{Contenido/Referencias}
%ANEXOS
\appendix
\renewcommand{\appendixname}{ANEXOS}
\chapter*{\appendixname}
\addcontentsline{toc}{chapter}{ANEXOS}
\markboth{ANEXOS}{ANEXOS}
\thispagestyle{empty}
ANEXO 1
\input{Contenido/Anexo Matriz de Consitencia}
Anexo 2
\input{Contenido/Anexo 2}



\end{document}
